\section{The inertia and monodromy conditions for the original rows}

% BEGIN I86A1
\paragraph{I86A1. The scalar on all inertia elements away from the coefficient prime.}
Fix a member, a coefficient index $i$ and a place $w\nmid\ell_i$.
Write $\xi_w$ for the local component of the tuple character at
$j=j(i)-1$ and $\lambda_w$ for its Weil character under the
member's local reciprocity convention. For every $s\in I_w$,
\[
                  \lambda_w(s)=\iota_i(\chi_1(i,s)).
\]
Indeed M84T10 identifies the global arithmetic line with the
$\ell_i$-adic avatar of the algebraic Hecke character obtained
from $\mu_C(i)\eta_T^{-1}$ with shift $-r/2$. Its proof compares
the specified split Frobenius values on a dense set; this part
uses only the coefficient index $i$, including when $\ell_i$
differs from both type primes. For an equivalent way to fix the
reciprocity convention, apply the rank-one Hecke-to-Galois
comparison to this algebraic character. The integrality of its
two infinity exponents follows from M84T6 and the computation
in M84T10. At a split good prime its Galois Frobenius value is
the unique Satake eigenvalue of that shifted character, exactly
the defining test for $\chi_1$. Density makes the two continuous
rank-one representations identical. Away from $\ell_i$ the
rank-one local comparison identifies the Weil character with
that of the shifted local Hecke character. The real shift and
$\eta_T^{-1}$ are unramified, and the native character reading
identifies $\mu_C(i)$ with the selected tuple. Restricting to
inertia gives the asserted equality with the original $\iota_i$
on the coefficient ring. Equivalently, in the idelic construction
the correction for the infinite exponents is supported above
$\ell_i$, so on units at $w$ only the same finite local
reciprocity value survives.
% END I86A1

% BEGIN I86A2
\paragraph{I86A2. Two spherical factors at the auxiliary places.}
Let $w\in S_i\cup Q_i$ satisfy the first disjunctive test in
$\mathrm{MonodromyInertiaType}$: either $w\in Q_i$, or
$\mathrm{Del}(i)$ holds and $w$ lies above $p_\delta$ or,
when $\ell_i=p_\delta$, above $p'_\delta$. Every such place
splits and is away from $\ell_i$. It is also away from both
type primes. For a $Q_i$ place this follows since all places over
a type prime away from $\ell_i$ belong to $S_i$ and $S_i,Q_i$
are disjoint. For the other disjunct it is the distinctness of
the auxiliary and type primes. Since $w\nmid\ell_i$, the
remaining $S_i$ case is over the active auxiliary prime, including
the switch to $p'_\delta$.

The unit character $\xi_w$ is tame and ramified here. In the
tuple $\mu_j^0\delta_{\rm act}\psi$ the first factor is
unramified away from the type primes. At a $Q_i$ place the
auxiliary factor is unramified by disjointness from $S_i$ and
$\psi$ is ramified by its exact $Q$-support. At a tested $S_i$
place $\psi$ is unramified and $\delta_{\rm act}$ has inertia
order 15. All these characters are tame on the indicated
units. Conjugate places carry inverse unit characters under
the CM identification, by the tuple's conjugate self-duality.
Write $k_w$ for the residue field. In a split coordinate supplied
by the level, the parahoric reducing to the $(r,t)$ parabolic
acts on any original level-character form by
$\theta(\det D)$, where $D$ is the lower diagonal block and
$\theta$ is the nontrivial character of $k_w^\times$ induced by
the tuple. Indeed the tested place belongs to the tuple's ramified
set by what we have just shown, so the member's $\mathrm{mt.para}$
applies at the place underneath it. Any local element of that
parahoric can be inserted as a single-place unitary adele. Its
principal congruence subgroup acts trivially. For every original
representative $P$ a nonzero vector in $\mathrm{typeV}(i,P)$
gives a nonzero vector with these transformations in its local
component $\pi$, by the isotypic projection argument of T86C3.

Here is the finite and local argument for any such residue field.
Its finite parabolic invariants contain
$1_{\mathrm{GL}_r(k_w)}\boxtimes(\theta\circ\det)$.
These finite factors have disjoint cuspidal supports, given by
$r$ trivial torus characters and $t$ copies of $\theta$.
The Mackey calculation in T86C5 applies over $k_w$, so their
finite parabolic induction occurs irreducibly in $\pi^{K_1}$.
In the multisegment support of $\pi$, each cuspidal factor has
principal congruence invariants, by exactness of compact
invariants and Iwasawa decomposition. For a depth-zero
supercuspidal, these invariants are its finite cuspidal type,
by depth-zero compact induction and the strict-gap Cartan
calculation in G85M2, valid for a local field with any residue
field. Finite parabolic induction of the support thus has the
same finite cuspidal support as its constituent just detected.
All factors here have rank one, on the two unramified-twist
lines of $1$ and $\xi_w$. This step uses no inertia-label or
large-residue-characteristic formula at the auxiliary prime.
Separate the lines by the unlinked-product theorem. The resulting
normalized product has blocks of ranks $r,t$. Exactness on the
maximal compact quotient and the identity-cell calculation in
T86C5 give the individual finite constituents $1$ and
$\theta\circ\det$ on those blocks. The first local representation
is therefore spherical and the second is spherical after twisting
by $\xi_w^{-1}\circ\det$. Both are supported on characters, so
the local correspondence's spherical, twist and product laws apply.
Their parameters have zero monodromy, the first is unramified
and the second acts on inertia through $t$ copies of $\lambda_w$.

If the chosen level coordinate corresponds to the other CM
place, use the inverse transpose between the two split coordinates
and tuple conjugate self-duality. By Gelfand--Kazhdan and the
duality and field-transport laws, as in T86C5, the dual parameters
give the two conclusions at $w$ with its own $\lambda_w$.
These laws apply here even for arbitrary residue degree since the
supports just obtained consist of characters. Transport by the
member's local reading and I86A1 give, for each $s\in I_w$,
\[
 \mathrm{mono0}(\mathrm{locV}(i,P,w)),\qquad
 \mathrm{inePol}(\mathrm{locV}(i,P,w),s)
       =(X-1)^r (X-\iota_i(\chi_1(i,s)))^t .
\]
% END I86A2

% BEGIN I86A3
\paragraph{I86A3. Extracting only the spherical block under the multiplicity guard.}
Let $w\in S_i$ lie above a type prime $g$. In particular
$w\nmid\ell_i$. Suppose the multiplicity of
$\alpha=\iota_i(\chi_1(i,\gamma))$ in the inertia polynomial
of $\mathrm{locV}(i,P,w)$ is at most $t$ at the designated
$\gamma=\mathrm{gam0}(i,w)$, with the native tame-generator
guard. We exhibit the two required local parameters at this
place. First work in the designated lift of the split pair.
The local component has a nonzero parabolic-radical fixed
vector of lower type $B=\theta\circ\det$, by the block argument
of T86C3 applied to the identity operator. Decompose under
$\mathrm{GL}_r(\mathbb F_g)$ and choose any irreducible top
type $A$ detected by this vector. For arbitrary $A$, finite
Frobenius reciprocity yields a common constituent of $A*B$ and
$\pi^{K_1}$. Its finite cuspidal support is the support of
$A$ plus $t$ copies of $\theta$. By the depth-zero support
calculation of G85M2, the number of factors on the local
unramified-twist line of $\xi$ equals the number of finite
cuspidal factors $\theta$ in this multiset. Each factor on that
line contributes one eigenvalue of $\lambda$ at \emph{every}
inertia element, since an unramified twist is trivial there;
the segment shifts preserve inertia as well. For the conjugate
coordinate the duality and transport argument of T86C5 gives
the same count with the conjugated inverse character, which is
$\lambda_w$ in that coordinate. All relevant support factors
are good here since the completion is $\mathbb Q_g$ and the
cuspidal factors have $K_1$-fixed vectors.

If $A$ contained even one extra $\theta$ in its cuspidal
support, the original inertia polynomial at $w$ would therefore
be divisible by $(X-\alpha)^{t+1}$, by I86A1. This contradicts
the assumed root multiplicity. Consequently the two finite
supports are disjoint and the second unramified-twist line has
exactly $t$ factors. The finite Mackey and unlinked-product
argument in T86C5 now gives a rank-$r$ first block and a rank-$t$
second block spherical after its character twist is removed.
No assertion about squarefreeness or monodromy of the first
block was required. Take the corresponding two Weil--Deligne
parameters, dualizing and transporting if $w$ is the conjugate
place. Their inverse images $a,b$ under the member's local
parameter identification satisfy
\[
 \begin{gathered}\mathrm{locV}(i,P,w)=a\oplus b,\quad
 \dim a=r,\quad\dim b=t,\quad N_b=0,\\
 \mathrm{inePol}(b,s)=(X-\iota_i(\chi_1(i,s)))^t\quad(s\in I_w).\end{gathered}
\]
In particular the proof uses the generator in the actual test;
equality with the Green-label generator is unnecessary.

Finally direct sums cancel at every place in the native parameter
space. Under its Weil--Deligne identification this is
Krull--Schmidt cancellation: the finite dimensional representations
with their monodromy operator form a category in which idempotents
split and each endomorphism has its invariant Fitting decomposition.
Finite dimension gives a finite decomposition into indecomposables,
whose endomorphism rings are local by Fitting's lemma, giving
uniqueness with multiplicities. Subtracting the multiplicities of
the common summand proves cancellation. I86A2, the displayed
decomposition and this unrestricted cancellation are the three
clauses of $\mathrm{At.MonodromyInertiaType}$. They prove the
original field $\mathrm{RowsB.a4}$ on every selected member.
% END I86A3
